\documentclass[a6paper,fontsize=10pt,twoside,openright]{scrbook}
\usepackage[a6paper,top=1cm,bottom=1.9cm,left=1cm,right=1cm,footskip=1cm,headsep=0.2cm]{geometry}% add showframe
\usepackage[swedish]{babel}
\usepackage{graphicx,tocloft,metalogo,titletoc,fontspec,titlesec,eso-pic,scrpage2,anyfontsize,microtype,array,marvosym,mdframed,textcomp,amssymb,enumitem}
\usepackage[columns=1]{idxlayout}

\newgeometry{top=1cm,bottom=1cm,left=1cm,right=1cm,footskip=0pt,headsep=0.2cm}


%\setlength{\emergencystretch}{2em}

\usepackage{letltxmacro}

\LetLtxMacro\origttfamily\ttfamily
\DeclareRobustCommand*{\ttfamily}{%
  \origttfamily
  \hyphenchar\font=`\-\relax
  \fontdimen3\font=.25em\relax
  \fontdimen4\font=.167em\relax
  \fontdimen7\font=.167em\relax
}

%% Font: Baskervald X
%% \usepackage[lf]{Baskervaldx} % lining figures
%% \usepackage[bigdelims,vvarbb]{newtxmath} % math italic letters from Nimbus Roman
%% \usepackage[cal=boondoxo]{mathalfa} % mathcal from STIX, unslanted a bit
%% \renewcommand*\oldstylenums[1]{\textosf{#1}}

%\usepackage[ebgaramond-maths]{mathdesign}
\usepackage[bitstream-charter]{mathdesign}

\usepackage[htt]{hyphenat}
\tolerance=1
\emergencystretch=\maxdimen
\hyphenpenalty
\exhyphenpenalty


\usepackage{setspace}
\renewcommand{\baselinestretch}{0.98} 

%\usepackage{showframe}
\usepackage{makeidx}
\makeindex

\let\oldemptyset\emptyset
\let\emptyset\varnothing

%% eso-pic for background on chapter pages %%
\newcommand\BackgroundPic{\includegraphics[width=\paperwidth,height=\paperheight,keepaspectratio]{elements/chapter2.pdf}}

%% scrpage2 pagestyling %%
\pagestyle{scrheadings}
\clearscrheadfoot

%% change empty page produced by \cleardoublepage into scrheading style %%
\makeatletter
\let\ps@empty\ps@scrheadings
\makeatother

\newcommand{\rpt}{
        \raisebox{.2ex}{:}
        \hspace*{-0.5ex}\raisebox{-.4ex}{\rule{.1ex}{2.3ex}\,\rule{.1ex}{2.3ex}}}

\newcommand{\revrpt}{
        \raisebox{-.4ex}{\rule{.1ex}{2.3ex}\,\rule{.1ex}{2.3ex}}
        \hspace*{-0.5ex}\raisebox{.2ex}{: }}

%% chapter and section styling %%
\titleformat{\chapter}[hang]{}{}{0pt}{\centering\fontsize{30}{70}\rmfamily\textbf}
\titleformat{\section}[hang]{\rmfamily\fontsize{13}{15}\bfseries}{}{0pt}{\textbf}
\titlespacing*{\section}{0pt}{0pt}{0pt}
\titlecontents{chapter}[0em]{\addvspace{3pt}}{\contentslabel{}}{}
              {\titlerule*[0.3pc]{.}\hspace*{-4pt}\contentspage}

%\newfontfamily\chapterfont[]{ChopinScript}

%\addtokomafont{pagenumber}{\oldstylenums}
              
%% Command for chapter page, will remove foot and head %%
\newcommand{\chapterpage}[2]{
  \newpage
  \null
  \cleardoublepage
  %\AddToShipoutPicture*{\BackgroundPic}
  \chapter[#1]{#2}
  \renewcommand{\leftmark}{#1}
  \newpage
  \ohead{\textnormal{\textsc{\scriptsize\leftmark}}}
  \ofoot[\pagemark]{\textsc{\scriptsize\pagemark}}
  %\cfoot{\centering\includegraphics[keepaspectratio]{elements/footelement2.pdf}}
}
\newcommand{\hiddenchapterpage}[1]{
  \newpage
  \null
  \cleardoublepage
  %\AddToShipoutPicture*{\BackgroundPic}
  \chapter*{#1}
  \renewcommand{\leftmark}{#1}
  \newpage
  \ohead{\textnormal{\textsc{\scriptsize\leftmark}}}
  \ofoot[\pagemark]{\textsc{\scriptsize\pagemark}}
  %\cfoot{\centering\includegraphics[keepaspectratio]{elements/footelement2.pdf}}
}
\makeatletter
\renewcommand\chapter{\if@openright\cleardoublepage\else\clearpage\fi
  \thispagestyle{empty}
  \global\@topnum\z@
  \@afterindentfalse
  \ofoot[\pagemark]{\textsc{\scriptsize\pagemark}}
  \cfoot{}
  \ohead{}
  \secdef\@chapter\@schapter}
\makeatother

%% Command for songs %%
\newcommand{\visa}[1]{
  \index{#1}
  \section{#1}
}
\newcommand{\fvisa}[2]{
  \index{#1}\index{#2}
  \section{#1}
}
\newcommand{\svisa}[3]{
  \index{#1}\index{#2}
  \section{#3}
}

%% TOC %%
\setcounter{tocdepth}{0}
\tocloftpagestyle{scrheadings}

%% MISC %%
\pagenumbering{arabic}
\renewcommand{\chaptername}{}
\renewcommand{\thechapter}{}
\renewcommand{\thesection}{}

\makeatletter
\renewenvironment{theindex}
                 {{\noindent\rmfamily{\fontsize{13}{15}\textbf{TITEL och FÖRSTARADEN REGISTER}}\vspace*{10pt}\par\textbf{\normalsize Symboler}\par}
                   \@mkboth{\MakeUppercase\indexname}
                           {\MakeUppercase\indexname}
                           \let\item\@idxitem}
                 {}
\makeatother

\newmdenv[
  topline=false,
  bottomline=false,
  rightline=false,
  linewidth=1pt,
  skipabove=5pt,
  skipbelow=3pt,
  innertopmargin=0pt,
  innerbottommargin=0pt
]{leftborder}
\makeatletter

\newcommand*\cleartoleftpage{%
  \clearpage
  \ifodd\value{page}\hbox{}\newpage\fi
}

\begin{document}
\setlength{\parindent}{15pt}
\vspace*{6.5cm}
\hspace*{0.9cm}
\includegraphics[keepaspectratio,width=7cm]{elements/logo.pdf}
\clearpage
\noindent Denna förträffliga oersättliga Manual tillhör:
\ohead{\textnormal{\textsc{\scriptsize\leftmark}}}
\ofoot[\pagemark]{\textsc{\scriptsize\pagemark}}
\clearpage
%\section{Datavetare är skapande människor}\vspace{10pt}
\noindent{\ttfamily Datavetare är skapande människor\\ \\ Datavetare
  skapar. De skapar nya program som världen aldrig sett tidigare, för
  att låta datorer göra saker som de aldrig gjort tidigare. De forskar
  fram nya kunskaper och öppnar nya fält för framtidens
  forskning. Datavetare skapar de verktyg som gör det möjligt att lösa
  nya problem. Nya programmeringsspråk, nya beräkningsverktyg, bättre
  kompilatorer, bättre databaser, nya operativsystem, bättre
  utvecklingsverktyg, distribuerade system... Listan kan göras hur
  lång som helst.\\ \\ Datavetenskapens handfasta mål är att förbättra
  den programvara som vi kör på våra datorer. Vad som krävs är att
  programmen gör vad de skall (korrekthet), att de inte gör fel om
  något oförutsett inträffar (pålitlighet), och att de arbetar
  effektivt (effektivitet). Dessutom skall programmen utvecklas på ett
  effektivt sett.}\\
\\
{\footnotesize\textit{Ur ``Vad är en datavetare egentligen''\\
    BlurGel nr 0, 1995}}
\clearpage
%\newgeometry{top=1cm,bottom=1.2cm,left=1cm,right=1cm,footskip=0pt,headsep=0.2cm}
\null
\vfill
    {\noindent\small\centering
      Redaktör: ?\\
      Grafisk formgivare: Per Bergqwist\\
      Logotyp av Jens Bohlin\\
      Producerad med \XeTeX\\
      v1.0: 1994-11-21 (1000 exemplar)\\
      v1.1: 1996-01-02 (1000 exemplar)\\
      v2.001: 2001-11-15 (1000 exemplar)\\
      v11111011111: 2015-??-?? (1200 exemplar)\\
      Tryckeri: Lyxo AB\par
    }
%\restoregeometry
\cleardoublepage
%\cfoot{\centering\includegraphics[keepaspectratio]{elements/footelement2.pdf}}
\section{FÖRORD}\vspace{10pt}
%\noindent{\large\textrm{\textbf{FÖRORD}}\vspace*{10pt}}\\
\hspace{\indent} ost
\newpage
\section{FÖRORD från v2.001}\vspace{10pt}
%\noindent{\large\textrm{\textbf{FÖRORD från v2.001}}\vspace*{10pt}}\\
\hspace{\indent}Allt sedan Manualen v1.1 kom till världen har det
funnits en skara människor som ägnat sina liv åt att vänta på dess
efterföljare. Denna skaras oförtrutliga tillväxt utmynnande i att ett
formaliserat uppdateringsarbete satte igång våren 1999. Den 4 december
-99 samlades en liten men representativ skara från C i Linköping, DVP
i Uppsala och C i Umeå för att samla alla tankar och idéer som
framförts på den då flitigt utnyttjade mailinglistan
\texttt{manual@acc.umu.se}. Efter detta möte trodde vi att arbetet
skulle leda till en färdig Manual våren 2000. Så blev det dock inte på
grund av en utdragen och högst osaklig debatt om en del låtar som
redan var bortplockade. Nu, över ett år senare, är vi som har
funderat, fipplat och till och med slitit(?) med denna nya bok
lyckliga över att få presentera Manualen version 2.001\\ \indent En
betydande samling sånger som Manualen kräver naturligtvis en avsevärd
mängd arbete att sammanställa, faktum är att antalet välvilliga
medhjälpare från alla berörda order förmodligen inte går att passa in
i en nibble. Det stora antalet inblandade parter gör att en uppräkning
är dömd att innehålla felaktigheter, dock gör vi nedan ett dj-ärvt
försök samtidigt som vi hoppas att alla vi glömt kan förlåta oss i
någon framtid.
\newpage
\section{FÖRORD från v1.0}\vspace{10pt}
%\noindent{\large\textrm{\textbf{FÖRORD från v1.0}}\vspace*{10pt}}\\
\hspace{\indent}Detta är en gemensam sångbok för samtliga datavetare
i Sverige; C i Lindköping, DVP i Uppsala och C i Umeå. Vi har genom
denna bok fått ytterligare en sak gemensamt.\\ \indent Att skapa en
gemensam sångbok har diskuterats flerga gånger, jag minns att
kontakter tog under 100-köret 1990 eller var det 1991?  Saker som
händer under 100-kör är inget man minns, utan saker som man senare
tror sig har hänt?! Den gången brakade arbetet samman efter några
veckor. Men idén levde kvar och gav 3 (eller 4) år senare detta som
resultat.\\ \indent Sångboken bygger på ``DVL:s Lilla blå v3.0\ss''
som aldrig han bli tryckt, innan arbetet kom igång för att skapa denna
bok. Dock skall sägas att denna bok är betydligt större än ``Lilla blå
v3.0\ss'', omkring 200 sidor för att inte vara exakt.\\ \indent När
\textit{manualen} trycks i december 1994 är boken den största sångbok
som finns runt om i vårt rike! Ok Psalmboken är större, men inte lika
roligt att ha på gasquer!\\ \indent Vi som varit ansvariga för denna
bok, har försökt att i så stor utsträckning som möjligt att få med
melodier, låtarnas ursprung, och hur man skall sjunga. Allt för att
underlätta för folk som inte är i kontakt med lika mycket sånger som
vissa av oss andra.\\ \indent Det är många som har hjälpt till med
denna bok, inte enbart datavetare utan även personer från helt andra
linjer och orter, jag vill rikta ett speciellt tack till Er.
\newpage
\indent Manualens generalgenomsjungning ägde rum på\\ Strålsnäs
herrgård den 18 November 1994, då vi avnjöt en underbar
rådjursmiddag. På denna sittning sjöng vi så mycket vi orkade och hade
efter ca 4 timmar traverserat ca 40 sånger.\\ \indent Med på denna
sångboksgenomsjungning och därav lika (o)ansvariga som jag till de
eventuella felen var,\\ \textit{Micke Liljeblad, Mats Hägglund, Jakob
  Engblom, Pär\\ Mattsson, Martin Davidson, Anders Jackson,
  Björn\\ Knutsson, Måns Engmem, Måns af Klercker, Jonas\\ Högström,
  Calle Englund, Andreas Gref} och sist jag själv.\\ \indent Innan jag
kommer till alla som skall tackas vill jag bara nämna att en sångbok
är till för att sjungas ur, detta betyder inte att man behöver
sjunga. Följande personer är jag evigt(?)  tacksam till;\\

\noindent
\begin{tabular}{@{}p{0.5\textwidth}p{0.5\textwidth}@{}}
  Anders Jansson & Anders Wilhelm\\
  Bertil Spolander & Björn Knutsson\\
  Björn Torkelsson & Erik Åhlin\\
  Göran Svensson & Håkan Lindholm\\
  Ingela Ljungkvist & Jan-Olof Flink\\
  Jens Bohlin & Joakim Thörnkvist\\
  Jonas Högström & Lars Hägglund\\
  Måns af Klercker & Olof Appelblad\\
  Pär Mattson
\end{tabular}
\newpage
\indent Dessutom vill jag tacka alla som lusläst boken, sammt alla Ni
som gjorde förhandsbeställningar. Det finns säkert många andra som jag
missat och jag är precis lika tacksam till er som alla
andra!\\ \indent Om ni har några bidrag eller saker som kan vara värda
att rätta till så skicka dessa till någon av festerierna nedan;\\

\noindent
\begin{tabular}{@{}p{0.3\textwidth}p{0.7\textwidth}@{}}
  Piraya klubben & piraya@ts.umu.se\\
  Klubbverket & dvl-klubbverket@student.docs.uu.se\\
  CC & cc@cyd.liu.se
\end{tabular}\\

\indent Till sist, kom ihåg att hela sångboken och fellista finns på:
\texttt{http://www.ts.umu.se/\raisebox{.4ex}\texttildelow manualen/}\\ \\ Umeå den 21 November
1994\\ \\ NeNNe
\cleartoleftpage
\section{TRAVERSERING}\vspace{10pt}
\hspace{\indent}Genom hela denna bok gäller följande traverseringsregler. I sånger där
en försångare eller dylikt skall sjunga ensam används fet text. Ifall
något speciellt såsom en handling inträffar vid ett ord, så används
(kursiv), där ordet för handling eller beteende står inom ( ). En
sådan handling gäller oftast maximalt raden ut. Ifall en längre
förklaring behövs kursiveras texten. Beskrivningen finns sedan efter
sångens slut. I denna beskrivning kan mycket annat matnyttigt finnas,
läs den gärna.\\ \indent I vissa sånger skall könen sjunga olika
rader, i dessa sånger används {\Large\Male} för manligt kön och där
kvinnligt kön skall sjunga används givetvis då {\Large\Female}.
\begin{leftborder}
  \hspace{\indent}Då det ibland kan förekomma lokala variantioner på vissa sånger, så
  används en  vertikal linje i högra kanten för att markera att detta
  är olika alternativ. Ibland betyder linjen att det är raderna som är
  alternativen, och ibland betyder linjen att det är olika
  versalternativ. Men den kloke läsaren förstår säkert när vad är vad.
\end{leftborder}
\hspace{15pt} För att göra allt ännu svårare så skall man oftast sjunga med stark
stämma där VERSALER används, men observera att man inte alltid skall
göra detta. Vem har sagt att det skall vara lätt?
\newpage
%% begin toc on odd page and style the title %%
\cleardoublepage
\renewcommand{\contentsname}{\vspace{-2.17cm}\rmfamily{\fontsize{13}{15}\textbf{INNEHÅLL}}\vspace{-1.2cm}}
\tableofcontents

%% Remove the toc title from heading %%
\renewcommand{\leftmark}{}

\setlength{\parindent}{0pt}
\raggedbottom
\chapterpage{Av tradition}{Av tradition}
\null
\vfill
\hfill\begin{minipage}[c]{6.7cm}
\textit{Vad bjuder oss uppriktigt Afrika?\\
Vad visa kan Amerika?\\
Vad Asien? Vad allt Europa?\\
Jag trotsar öppet alltihopa.\\
Men Skandinavien - det är alladar!\\
Blott Sverige svenska krusbär har.\\
\\
Carl Love Almquist}
\end{minipage}\hfill
\vfill
\null
\newpage
\visa{DU GAMLA, DU FRIA}
Du gamla, Du fria, Du fjällhöga nord\\
Du tysta, Du glädjerika sköna!\\
Jag hälsar Dig, vänaste land uppå jord,\\
||: Din sol, Din himmel, Dina ängder gröna. :||\\
\\
Du tronar på minnen från fornstora dar,\\
då ärat Ditt namn flög över jorden.\\
Jag vet att Du är och Du blir vad du var.\\
||: Ja, jag vill leva jag vill dö i Norden. :||\\
\\
Jag städs vill dig tjäna mitt älskade land,\\
din trohet till döden vill jag svära.\\
Din rätt, skall jag värna, med håg och med hand,\\
||: din fana, högt den bragderika bära. :||\\
\\
Med Gud skall jag kämpa, för hem och för härd,\\
för Sverige, den kära fosterjorden.\\
Jag byter Dig ej, mot allt i en värld\\
||: Nej, jag vill leva jag vill dö i Norden. :||\\
\\
{\footnotesize\textit{Text: Richard Dybeck, 1844\\
\\
Sveriges nationalsång av tradition sedan 1866.\\
De två sista verserna sjungs sällan.}}
\clearpage

\newpage
%\fvisa{SVERIGES FLAGGA}{Flamma stolt mot dunkla skyar}
\vspace{10pt}
Flamma stolt mot dunkla skyar\\
lik en glimt av sommarens sol!\\
Över Sveriges skogar, berg och byar,\\
över vatten och viol!\\
Du som sjunger, när Du bredes\\
som vår gamla lyckas tolk.\\
Solen lyser! Solen lyser!\\
Ingen vredes åska slog vårt tappra folk!\par
\vspace{10pt}
Flamma högt vårt kärlekstecken!\\
Värm oss, när det blåser kallt!\\
Stråla ut de blåa vecken\\
kärlek mera stark än allt!\\
Sveriges flagga! Sveriges ära!\\
Fornklenod och framtidstolk!\\
Gud är med oss! Gud är med oss!\\
Han skall bära stark vårt fria svenska folk\par
\vspace{10pt}
{\footnotesize\textit{Text: K.G. Ossiannilsson\\
Musik: Hugo Alfvén}}

%\fvisa{JÄMTLANDSSÅNGEN}{Så tåga vi tillsammans bort}
\vspace{10pt}
Så tåga vi tillsammans bort\\
mellan Jämtlands gröna ängar\\
bort mellan nyland som prunka\\
fulla av bröllopsblomsters prakt.\\
Så skådom nu med gamman hän\\
emot berg i blåa fjärran hän,\\
över sjöar, strömmar, skogar\\
jämt kring bygder på vakt.\\
\\
\revrpt Fagert är landet\\
som blev vår lott och arvedel.\\
Så firom dess fägring nu\\
med sång och stråkars spel.\\
Så tändom ånyo det hopp\\
som våra fäder närt.\\
För slit och mödor\\
av fröjd och sol\\
ett mått oss beskärt\rpt\\
\\
{\footnotesize\textit{Text \& Musik: Wilhelm Peterson-Berger}}
\clearpage

%\input{songs/jamtlandssången.tex}
\input{songs/studentsången.tex}
\newpage
\visa{UNDER SVEA BANÈR}
\vspace{10pt}
Under Svea Banér Himlen seger oss ger;\\
Då för Konung och Land Äran lyfter sin hand.\\
Ännu Svearnes mod sif bereder en Stod\\
Utaf Lagrarne höjd, fast besglad mes blod.\\
Himlen gifve oss frid!\\
Men om den icke vinne utan vapen och strid,\\
Blifve Segern då vår, ögat offre sin tår\\
Den som faller, ock rätta till vår tacksamhet får.\\
\\
{\footnotesize\textit{Text: Samuel Ödmann\\
Musik: Johann Christian Friedrich Haeffner\\
\\
Anses vara den första upsaloensiska studentsången. Musiken hade
skrivits tidigare som soldatkör till Hæffners opera “Renaud” (\oldstylenums{1801}).\\
\\
Framfördes första gången, mitt under kriget mot Ryssland,
den \oldstylenums24} oktober \oldstylenums{1808}. Detta skedde vid en
uppvakting för fältmarskalkten Klingspor på genomresa från Finland
till Stockholm som gästade landshövding Uppsala slott.  Uppsalas
studenter marscherade då upp till slottet. Sångerna gick i täten och
sjöng denna sång.}}

\newpage
\visa{O, GAMLA KLANG- OCH JUBELTID}
{\footnotesize\textit{Melodi: Oh alte Burschenherrlichkeit}}\par
\vspace{10pt}
O gamla klang och jubeltid\\
ditt minne skall förbliva\\
och än åt livets bistra strid,\\
ett rosigt skimmer giva.\\
Snart tystnar allt vårt yra skämt,\\
vår sång blir stum, vårt glam förstämt.\\
O, jerum, jerum, jerum.\\
O, quae mutatio rerum!\par
\vspace{10pt}
Var äro de som kunde allt,\\
blott ej sin ära svika.\\
Som voro män av äkta halt\\
och världens herrar lika?\\
De drogo bort från vin och sång\\
till vardagslivets tråk och tvång.\\
O, jerum, jerum, jerum.\\
O, quae mutatio rerum!\\
\begin{leftborder}
\begin{tabular}{l l}
  (Filosofer)  & Den ene vetenskap och vett\\		
               & in i scholares mänger,\\
  (Jurister)   & Den andre i sitt anlets svett\\
               & på paragrafer vränger,\\
  (Teologer)   & en plåstrar själen, som är skral,\\
  (Medicinare) & en lappar hop dess trasiga fodral;\\
  (Alla)       & O, jerum, jerum, jerum,\\
               & O, quae mutatio rerum!
\end{tabular}
\end{leftborder}
\newpage
Men hjärtat i en sann student,\\
kan ingen tid förfrysa.\\
Den glädjeeld, som där har tänt,\\
hans hela liv skall lysa.\\
Det gamla skalet brustit har\\
men \textit{kärnan} finnes frisk dock kvar,\\
och vad han än må mista,\\
den skall dock aldrig brista!\par
\vspace{10pt}
Så sluten, bröder, fast vår krets,\\
till glädjens värn och ära!\\
Trots allt vi tryggt och väl tillfreds,\\
vår vänskap trohet svära.\\
Lyft bägar'n högt, och klinga vän!\\
De gamla gudar leva än\\
\revrpt bland skålar och pokaler\rpt\par
\vspace{10pt}
{\footnotesize\textit{Text: August Lindh, 1925\\ Musik:
Eugen Höfling, 1825}}\par
\vspace{10pt}
{\footnotesize\textit{Vid ``kärnan' dunkas näven i bordet en
gång. Den sista versen sjungs stående samt stående på stolen, efter
sista versen sjungits sätter sig ingen åter till bords.}}

\chapterpage{Datavetarvisor}{Datavetarvisor}
\newpage
\thispagestyle{empty}
\null
\vfill
{\noindent\textit{\centering Fader Dator, som är i Salen,\\ helgad
    vare Din skärm,\\ tillkomme Ditt tangentbord,\\ ske Din vilja,
    så som i minnet,\\ så ock på nätet.\\ Vår dagliga utskrift giv
    oss idag\\ och förlåt oss våra misstag\\ trots att vi ej
    förlåta\\ dem som skyldiga äro.\\ Låt oss icke behöva
    vänta\\ och fräls oss ifrån dumpar.\\ Ty Salen är din\\ och
    Makten hos Handledaren\\ i Evighet.\\ UNIX.\par}}
\vfill
\null
\newpage
\fvisa{DATAVETARNAS HÄRJARVISA}{Nu ska vi kompilera}
      {\footnotesize\textit{Melodi: Gärdebylåten}}\\
\\
Nu ska vi kompilera,\\
supa och penetrera\\
söka rekusioner ända in\\
till fixpunktssemantik.\\
Av induktion och basfall\\
blir man med lätthet asknall.\\
Hashtabeller parsas rekursivt\\
ifrån hår till häl\\
\\
{\footnotesize\textit{Text: Pär Mattsson och Magnus Ingelbo från DVL Uppsala.}}

\vspace{15pt}
\fvisa{JUBILEUMSVISA (MADHOUSE)}{En datavetare, han bor vid FooBar}\par
{\footnotesize\textit{Melodi: En sockerbagare}}\par
\vspace{10pt}
{En datavetare, han bor vid FooBar\\
sitter och kodar, vill aldrig blir klar\\
har overallen, han har långt hår.\\
har redan suttit i många år\\
Och denna dagen den skall nu firas\\
hopp in i duschen, sig själv omsvidas\\
dags att klippa sitt långa hår\\
tänk att det redan gått tretti år\par
\vspace{10pt}
{\footnotesize\textit{Text: Erika Elmquist, Kristiina Ausmees och Jens
Rosén från DV Uppsala}}\par
{\footnotesize\textit{Skriven inför DV
Uppsalas Jubileumsgasque (30 år jubiléet)}}\\
{\footnotesize\textit{Sjungs med fördel med
amerikansk brytning}}

\newpage
\svisa{THE BASIC SONG@\texttt{THE BASIC SONG}}{10 LET oss nu fatta i våra glas@\texttt{10 LET oss nu fatta i våra glas}}{\texttt{THE BASIC SONG}}
{\footnotesize\textit{Melodi: Mors lilla Olle}}\\
\\
\texttt{10 LET oss nu fatta i våra glas}\\
\texttt{20 INPUT en klunk utav det som där has}\\
\texttt{30 IF du fått nog THEN till 50 min vän}\\
\texttt{40 ELSE GOTO-baka till 10 igen}\\
\texttt{50 END}\\
\\
{\footnotesize\textit{Inledande radnummer sjungs ej.\\ I övrigt följs
    kommandon, ex. drick vid ``\texttt{INPUT en klunk}...''}}

\vspace{15pt}
\svisa{GE MIG MERA"!}{Stackars kursexaminatör}{GE MIG MERA!}
{\footnotesize\textit{Melodi: Fula gubbar}}\par
\vspace{10pt}
Stackars kursexaminatör,\\
som skall få smaka på allt som jag gör.\\
Det blir Pascal, Prolog och Lisp och C.\\
En inlupp ska du se!\\
Jag kallas hacker, freak och nörd\\
och jag kallas även störd.\\
För mig så dansar UNIX balett.\\
Jag märker ej att jag luktar svett.\\
Med en turbo-hackad fönstermiljö\\
som en manlig laddad mö.\par
\vspace{10pt}
Ge mig mera! Kompilera!\\
Runt på Polacksbacken springer jag.\\
Snart kan jag avancera.\\
Datavetare jag blir en dag.\par
\newpage
På HTML är jag autodidakt.\\
Använder News för att få kontakt.\\
Jag låter GIF-bilder bli en ångventil,\\
för lust så pubertil.\\
Och när jag en gång kommer ut,\\
blir min ingångslön till slut,\\
så hög att till och med jag får en chans.\\
Amor vill ge mig en segerkrans,\\
när jag får med mig upp på mitt rum,\\
en ny Pentium.\par
\vspace{10pt}
Ge mig mera...\par
\vspace{10pt}
{\footnotesize\textit{Text: Måns af Klercker\\Framförd på DVP
      Uppsalas Jubileumsgask (15 år jubiléet)}}

\newpage
\fvisa{DATAVISA}{Jag är helröd och helt OK}
{\footnotesize\textit{Melodi: The lumberjack song}}\par
\vspace{10pt}
Jag är helröd och är helt OK.\\
Jag jobbar hårt och jag roar mig.\par
\vspace{10pt}
Han är helröd och är helt OK.\\
Han jobbar hårt och han roar sig.\par
\vspace{10pt}
Data är ball,\\
jag kan Pascal.\\
Till FooRum vill jag gå.\\
Där träffas alla vänner\\
som är från Uppsala.\par
\vspace{10pt}
Data är ball,\\
han kan Pascal.\\
Till FooRum vill han gå.\\
Där träffas alla vänner\\
som är från Uppsala.\par
\vspace{10pt}
För han är helröd och är helt OK.\\
Han jobbar hårt och han roar sig.\par
\vspace{10pt}
Min mattebok\\
den gör mig klok.\\
Jag läser semantik.\\
Jag går på föreläsning\\
och älskar juridik.\\
\newpage
Hans mattebok\\
den gör han klok.\\
Han läser semantik\\
Han går på föreläsning\\
och älskar (\textit{förvånat}) juridik?\par
\vspace{10pt}
Men han är helröd och är helt OK.\\
Han jobbar hårt och han roar sig.\par
\vspace{10pt}
Som ekonom\\
jag blir fantom.\\
Konkurser gör mig säll.\\
Till flickor blankt jag nekar,\\
jag älskar en tabell.\par
\vspace{10pt}
(\textit{förvånad}) Som ekonom\\
han blir fantom (\textit{ursinnigt}) näääh, buuuuh!\par
\vspace{10pt}
(\textit{hurtigt}) \revrpt Men han är helröd och är helt OK.\\
Han jobbar hårt och han roar sig.\rpt\par
\vspace{10pt}
{\footnotesize\textit{Modifierad av NeNNe, dvl89, inför
    DVL Umeås 5 - års jubileum.\\ Orginalet är en
    teknologvisa från Lund.\\ (Ljusblå, KS och Umeå) har bytas mot
    (helröd, FooRum och Uppsala) för att översättas till uppländska.}}

\newpage
\fvisa{WRITE IN C}{When I find my code in tons of trouble}
{\footnotesize\textit{Melodi: Let it be}}\par
\vspace{10pt}
When I find my code in tons of trouble,\\
friends and colleagues come to me,\\
speaking words of wisdom:\\
``Write in C.''\par
\vspace{10pt}
As the deadline fast approaches,\\
and bugs are all that I can see,\\
somewhere, someone whispers:\\
``Write in C.''\par
\vspace{10pt}
Write in C, write in C,\\
write in C, oh, write in C.\\
LOGO's dead and buried,\\
write in C.\par
\vspace{10pt}
I used to write a lot of FORTRAN,\\
for science it worked flawlessly.\\
try using it for graphics!\\
Write in C.\par
\vspace{10pt}
If you've just spent nearly 30 hours\\
debugging some assembly,\\
soon you will be glad to\\
write in C.\par
\newpage
Write in C, write in C,\\
write in C, yeah, write in C.\\
Only wimps use BASIC.\\
write in C.\par
\vspace{10pt}
Write in C, write in C \\
write in C, oh, write in C.\\
Pascal won't quite cut it.\\
write in C.\par
\vspace{10pt}
Write in C, write in C,\\
write in C, yeah, write in C.\\
Don't even mention COBOL.\\
Write in C.\par
\vspace{10pt}
{\footnotesize\textit{Text: Kriston J. Rehberg}}

\newpage
\fvisa{OM EMACS}{Emacs är en stor koloss}
{\footnotesize\textit{Melodi: Rullan går}}\\
\\
Emacs är en stor koloss.\\
Tugga på, tugga på.\\
Verkar aldrig komma loss.\\
Tugga på, tugga på.\\
För att flytta på ett tecken\\
går två meg i garbage-säcken.\\
Ner i skärmens undre ran’\\
står: garbage collecting, done!\\
\\
Väntar på’n till min pension.\\
Tugga på, tugga på.\\
Emacs gör en Sparcstation\\
Tugga på, tugga på.\\
snabb som en ABC-80.\\
Högprestanda blir så smått i\\
Emacs, den ser säkert till\\
ingen dödtid blir förspilld!\\
\\
Emacs skriven är i Lisp.\\
Tugga på, tugga på.\\
Inget rimmar alls på Lisp.\\
Tugga på, tugga på.\\
Om du in i kärnan trevar\\
finn en gnu som motorn vevar.\\
Den drar Emacs runt för hand!\\
Klart att det går trögt ibland!\\
\\
{\footnotesize\textit{Text: Ingemar Ragnemalm \oldstylenums{1992}}}

\newpage
\fvisa{1515}{Höstvindar vina, ursh vilken pina}
{\footnotesize\textit{Melodi: Vårvindar friska}}\\
\\
Höstvindar vina, ursh vilken pina,\\
jag tror jag stannar hemma idag.\\
Vid terminalen, inne i salen,\\
uppe i 1515.\\
Hackar på SPARC:en, skriver i C,\\
minneshantering, vad är nu det?\\
Trycker ctrl-x, trycker ctrl-s,\\
dags då att kompilera.\\
\\
gcc gnäller, lint bara skäller,\\
“Code way too ugly, won’t even try!”\\
Ändrar ett tal här, ändrar en rad där,\\
cast’ar en liten int.\\
Prova att kompilera igen,\\
titta det gick igenom ... nej men...\\
Oj! Nej! Vad händer? Gnistrar och bränner,\\
ajdå, vi fick en core-dump.\\
\\
{\footnotesize\textit{Text: Jesper Wilhelmsson, d97\\
    \\
    ... Ja, jag vet att det ska vara C-x...}}

\newpage
\fvisa{KOMMANDOTOLKEN}{Far, jag kan inte få upp min kommandotolk}
{\footnotesize\textit{Melodi: Far, jag kan inte få upp min kokosnöt}}\\
\\
Far, jag kan inte få upp min kommandotolk.\\
Alla sätt jag prövat har vart fel;\\
jag har tryckt på knappar, klickat på varje flik,\\
och jag har sagt, “kommandotolk, tack!”,\\
men skärmen är sig lik.\\
\\
Far, jag kan inte få upp min kommandotolk,\\
inte ens med länk till satellit.\\
På varje antenn, har jag satt kopplingssladdar, men,\\
av kommandotolken min syns inte ett skit!\\
\\
Av kommandotolken min syns inte ett ski-i-it!\\
Jag ser bara Emacs, Xfig och Mosaic.\\
Säg, minns du prompten, far?\\
Glöm bort den, är du rar.\\
för av kommandotolken min syns inte ett skit!\\
\\
Far, jag kan inte få upp min kommandotolk,\\
inte ens med nitroglycerin.\\
När jag hade laddat, mycket omsorgsfullt och väl,\\
bortåt jag sprang, sen sa’ det PANG,\\
och Polacksbacken brann ner...\\
\\
Far, jag kan inte få upp min kommandotolk,\\
för jag har ingen hand med dynamit.\\
Av vårt stolta hus, det återstår bara grus,\\
och av kommandotolken min syns inte ett skit!\\
\\
Av kommandotolken min syns inte ett ski-i-it!\\
Allt det andra har blitt lika indefinit (så att säga).\\
Internet står still, vad ska jag ta mig till?\\
Av kommandotolken min syns inte ett skit,\\
ett skit,\\
ett skit,\\
ett skit,\\
Av kommandotolken min syns inte ett ski-i-it!\\
\\
{\footnotesize\textit{Text: Johan Runeson, Nollespexet DVP Uppsala,
    \oldstylenums{1996}\\ \\ (Personligen tycker jag att det ska vara
    ’Fan, jag kan inte...’ /red. anm.)}}

\vspace{15pt}
\visa{MORS LILLA DATOR}
{\footnotesize\textit{Melodi: Mors lilla Olle}}\par
\vspace{10pt}
Mors lilla dator åt skogen gick,\\
mitt i programmet så sade det klick\\
svart bidde skärmen och minnet försvann,\\
den informationen kan ingen få fram.\par
\vspace{10pt}
Brummeli-brum vad brummar där?\\
Det sprakar och gnistrar, ett jordfel det är!\\
Blixtarna blå ifrån burken det slå,\\
synd att jag nu här ensammen stå.\par
\vspace{10pt}
Hyscheli-hysch vad prasslar här?\\
Fram väller pappret ur printern där!\\
Den har fått nippran av tecken så små,\\
jag tror att jag snart hemåt skall gå.

\newpage
\fvisa{ANOTHER GLITCH IN THE CALL}{We don't need no indirection}
{\footnotesize\textit{(sung to the tune of a not so recent Pink Floyd song)}}\par
\vspace{10pt}
We don't need no indirection\\
We don't need no flow control\\
No data typing or declarations\\
Hey! Did you leave the lists alone?\par
\vspace{10pt}
All in all, it's just a pure-LISP\\
function call\par
\vspace{10pt}
We don't need no side effect-ing\\
We don't need no scope control\\
No global variables for execution\\
Hey! Did you leave those args alone?\par
\vspace{10pt}
All in all, it's just a pure-LISP\\
function call\par
\vspace{10pt}
We don't need no allocation\\
We don't need no special nodes\\
No dark bit-flipping in the functions\\
Hey! Did you leave the bits alone?\par
\vspace{10pt}
All in all, it's just a pure-LISP\\
function call
\newpage
We don't need no compilation\\
We don't need no load control\\
No link edit for external bindings\\
Hey! Did you leave that source alone?\par
\vspace{10pt}
All in all, it's just a pure-LISP\\
function call\par
\vspace{10pt}
{\footnotesize\textit{Text: DVLs tidning \#11subscript:8}}

\vspace{15pt}
\fvisa{TUPPENS KLAGOVISA}{Tuppens väckarklocka ringer sex varje morgon}
{\footnotesize\textit{Melodi: Hoppsan vilken dag}}\\
\\
Tuppens väckarklocka ringer sex varje morgon\\
då måste han gå upp till ännu en jävlig dag.\\
Han sätter sig på Teknikum och får genast ångest\\
rysningar i kroppen alla världens obehag.\\
\\
Ännu en föreläsning om fysik och sånt\\
men han tänker på en svettig datasal.\\
Tänk att sitta vid en dator i 1411\\
men miniräknarn är hans enda val.\\
\\
Tuppens verklighet det är ett liv utan glädje\\
gula overaller är ju inget man vill ha.\\
Jag borde ändå gjort som min gamla mamma sa,\\
en DVP:are é va man borde va\\
\revrpt en DVP:are é va man borde va\rpt

\newpage
\fvisa{HALLANDS-SHOT}{Ingen vet det jag vet en hemlighet}
{\footnotesize\textit{Melodi: Margareta}}\par
\vspace{10pt}
Ingen vet det jag vet en hemlighet\\
vad det var i glaset mitt som jag svept\\
kan inte hjälpa att jag känner det så här\par
\vspace{10pt}
Och den kick som jag fick gav mig allt\\
världen börjar snurra runt vad jätteballt\\
tänk att allting känns så härligt när\par
\vspace{10pt}
man dricker Rhoca-Gil ska du veta \\
då blir alla öppningar täta\\
står här bak' en dörr\\
knip' som aldrig förr\\
och kan inte pissa\par
\vspace{10pt}
pulsarna de bränner så heta\\
dricker du som jag ska du veta\\
suparna du tar\\
har du evigt kvar - inom dig\par
\vspace{10pt}
Så jag går ut igen till nästa bar\\
ännu en Hallands-shot jag genast tar\\
tänk att allting känns så härligt när\par
\vspace{10pt}
man dricker Rhoca-Gil ska du veta \\
då blir alla öppningar täta\\
står här bak' en dörr\\
knip' som aldrig förr\\
och kan inte pissa

pulsarna de bränner så heta\\
dricker du som jag ska du veta\\
suparna du tar \\
har du evigt kvar - inom dig\par
\vspace{10pt}
{\footnotesize\textit{Text: Anders Ramsell, Johan Runeson, Henrik Eriksson}}

\vspace{15pt}
\visa{DET VAR FEST NER I FOORUM}
{\footnotesize\textit{Melodi: Det var dans bort i vägen}}\par
\vspace{10pt}
Det var fest ner i FooRum på tisdagsnatten.\\
Upp i trappan gick låten av sången och skratten.\\
Det var tjo! Det var tjim! Det var hej!\\
Där satt dataloger och skåla och söp\\
och skråla och skräna tills hemåt de kröp\\
för dudeli dudeli dej. HEJ!\par
\vspace{10pt}
Där var Starka å Åhus å Renat å Bäsk,\\
där var Herrgårds och Pors som ej luktar av mäsk\\
där var OP och lite Mai Tai.\\
Där var Årsta och Vinbärs och liknande klägg,\\
där var Punsch och likör utav lingon och ägg\\
för dudeli dudeli dej. SKÅL!\par
\vspace{10pt}
{\footnotesize\textit{Text: Andreas Gref, dvl90 Uppsala}}

\newpage
\fvisa{DET SKA VA' HEMBRÄNT}{Jag har festat mycket, det har inte varit lätt}
{\footnotesize\textit{Melodi: Husvagn}}\par
\vspace{10pt}
Jag har festat mycket, det har inte varit lätt.\\
Blanda vin och whisky, det kan inte vara rätt.\\
Jag har festat på de allra konstigaste sätt,\\
men äntligen jag funnit hur man ska partajja rätt:\par
\vspace{10pt}
Det ska va hembränt, och supa till i några dar.\\
Det ska va hembränt, drick upp det sista ni har kvar.\\
Det ska va hembränt, och svina runt i overall.\\
Det ska va hembränt, drick nu vad ni tål!\par
\vspace{10pt}
I flera år så var vi lite fina där i kanten.\\
Vi smuttade på vinerna men spara aldrig panten.\\
Men sen vi lärde känna våran vackra overall,\\
så super vi nu landet runt på hemgjord alkohol.\par
\vspace{10pt}
Det ska va hembränt ...\par
\vspace{10pt}
Fem små groggar först för att komma igång med festen.\\
Fem små groggar till för att supa ikapp resten.\\
Fem små groggar sen för att kalla det en kväll.\\
Fem små groggar till så sitter vi i fyllecell.\par
\vspace{10pt}
Det ska va hembränt ...\par
\vspace{10pt}
{\footnotesize\textit{Text: Jesper Wilhelmsson, d97}}

\newpage
\fvisa{DATAVETARNAS PARADMARCH}{Ja, här kommer datavetare}
{\footnotesize\textit{Melodi: Fritjof Anderssons paradmarch}}\par
\vspace{10pt}
Ja, här kommer datavetare\\
med märken på vår dräkt.\\
Vi knappar kod, \\
vi hackar LISP\\
och surfar runt på nätet.\par
\vspace{10pt}
Vi lämnar pöbeln stående\\
och gapande förskräckt.\\
Det vi inte kan, \\
det kan ingen ann -\\
nu går vi upp i labbet.\par
\vspace{10pt}
Och när vi loggat ut,\\
ja då går vi på någon fest\\
och tar oss några bira,\\
för det gillar vi ju bäst!\par
\vspace{10pt}
För där är där vi inte är\\
och här är där vi är!\\
Vi knappar kod,\\
vi hackar LISP\\
och festar runt Sverige.\par
\vspace{10pt}
{\footnotesize\textit{Text: Selander, C93 \& Daniel ``Holm''
    Holmgren, C96.\\ Skrevs i baksätet på en Toyota Camry, på väg
    till linjekonferensen -97 i Uppsala. Omdöpt av Jubileumsgruppen
    inför 30 år jubiléet i Uppsala}}

\newpage
\fvisa{IT IS FREE}{When I find myself in front of Windows}
{\footnotesize\textit{Melodi: Let It Be}}\par
\vspace{10pt}
When I find myself in front of Windows,\\
Father Torvalds comes to me,\\
Speaking words of Linux, ``It is free!''.\\
And when hackers crack my system\\
through some strange NT obscurity,\\
he gives me the source-code, ``It is free!''.\par
\vspace{10pt}
It is free, it is free, it is free, it is free.\\
Linux is the answer, and it's free.\par
\vspace{10pt}
When the broken system crashes\\
and blue screens of death I see.\\
There is one true rescue, and it's free.\\
For though there may be doubts first\\
it will sure bring life to your PC\\
`Cause Linux is the answer, and it's free.\par
\vspace{10pt}
It is free, it is free, it is free, it is free,\\
You even get the source-code, it is free.\par
\vspace{10pt}
In the darkest night, when hacking,\\
Xfree stares right back at me,\\
Now I've installed Linux, I feel free.\\
No longer mad with my computer,\\
No more does it crash on me\\
Now I've installed Linux, it is free.
\newpage
It is free, it is free, it is free, it is free,\\
Open-sourced true power, it is free.\par
\vspace{10pt}
{\footnotesize\textit{Text: David Weinehall}}

\vspace{15pt}
\fvisa{OH, LJUVA ÖL}{Jag har nu fått en och annan öl}
{\footnotesize\textit{Melodi: Kung Louis sång}}\\
\\
Jag har nu fått en och annan öl,\\
som jag genast druckit har.\\
Ja, jag har druckit mycket jag.\\
Jag ingen droppe spar.\\
För nu vill jag bara festa,\\
tills jag far ikull.\\
Jag vill ej längre nykter va.\\
Jag vill bara vara full.\\
\\
Oh, ljuva öl, vad jag är glad åt öl.\\
För jag vill dricka den, och en till, öl.\\
Det vill jag nu, ett djur som jag,\\
det lär sig bra bli en fyllehund.\\
\\
Försök inte hindra mig gosse,\\
för nu så är det fest.\\
Och då man drycker dricka ska,\\
som gjorda är på jäst.\\
Så ge mig nu inget annat,\\
att hälla i mitt glas.\\
för nu så ska jag dricka ur,\\
och bli full som ett as.

\chapterpage{Snapsvisor}{Snapsvisor}
\section{SUPARNAS ORDNING}\index{SUPARNAS ORDNING}
\vspace{10pt}
\begin{enumerate}[leftmargin=0.65cm,topsep=0pt,itemsep=0pt,partopsep=0pt,parsep=3pt]
\item Helan
\item Halvan
\item Tersen.
\item Kvarten
\item Kvinten
\item Sexten
\item Septen
\item Rivan
\item Rafflan.
\item Rännan
\item Smuttan
\item Smuttans unge
\item Femton droppar
\item Lilla Manasse
\item Lilla Manasses Broder
\item Klämtaren
\item Kreaturens återuppståndelse
\item Den bleka dödens dryck
\item Krokofanternas intåg
\item Ett evigt liv
\end{enumerate}
\newpage
\section{ANVISNINGAR FÖR SNAPSENS TAGANDE}\index{ANVISNINGAR FÖR SNAPSENS TAGANDE}
\vspace{10pt}
\begin{itemize}[leftmargin=0cm,topsep=0pt,itemsep=2pt,partopsep=0pt,parsep=3pt,label={}]
\item Snapsen tages till fisk.
\item All mat är fisk, förutom ölsupa och pannkaka.
\item Den som äter ölsupa eller pannkaka, sup på grannens fisk.
\item Skulle grannen förtära ölsupa eller pannkaka, må även dessa rätter betraktas som fisk.
\item Desserten ätes med gaffel, så länge det går, därefter med sked.
\item Man må icke dricka ur sköljkopparna.
\item Ty den goda törsten bör släckas med ädlare fludium än tvättvatten.
\item Vatten betraktas aldrig som fisk.
\end{itemize}
\chapterpage{Ölvisor}{Ölvisor}

\chapterpage{Vinvisor}{Vinvisor}

\chapterpage{Avecvisor}{Avecvisor}

\chapterpage{Gasquevisor}{Gasquevisor}
\fvisa{HÄRJAREVISAN}{Liksom våra fäder, vikingarna i Norden}
{\footnotesize\textit{Melodi: Gärdebylåten}}\par
\vspace{10pt}
Liksom våra fäder, vikingarna i Norden.\\
Drar vi riket runt och super oss under borden.\\
Brännvinet har blitt ett elixir,\\
för kropp såväl som själ.\\
Känner du dig liten och ynklig på jorden,\\
växer du med supen och blir stor ut i orden.\\
Slå dig för ditt håriga bröst och bli \\
en man från hår till häl.\par
\vspace{10pt}
Ja, nu skall vi ut och härja,\\
supa och slåss och svärja,\\
bränna röda stugor, slå små barn och säga fula ord.\\
Med blod skall vi stäppen färga.\\
Nu äntligen lär ja'\\
kunna dra någon riktig nytta\\
av min Hermodskurs i mord.\par
\vspace{10pt}
Hurra, nu ska man äntligen få röra benen,\\
hela stammen jublar och det spritter i grenen.\\
Tänk att än en gång få spränga fram\\
på Brunte i galopp!\\
Din doft o kära Brunte är trots sin brist i hygienen,\\
för en vild mongol minst lika ljuv som syrenen.\\
Tänk att på din rygg få rida runt\\
i stan och spela topp.\par
\vspace{10pt}
Ja, nu skall vi ut och härja...\par
\vspace{10pt}
Ja, mordbränder är klämmiga ta fram fotogenen\\
och eftersläckningen tillhör just de fenomenen\\
inom brandmansyrket som jag tycker\\
det är nån nytta med.\\
Jag målar för mitt inre upp den härliga scenen;\\
Blodrött mitt i brandgult, ej ens prins Eugen en\\
lika mustig vy kan måla, \\
ens om han målade med sked.\par
\vspace{10pt}
Ja, nu skall vi ut och härja...\par
\vspace{10pt}
{\footnotesize\textit{Text: Hans Alfredsson och Levin\\ Ur Lundaspexet
    Djingis Kahn 1954, första versen ej ur originalet.}}\\
\index{DATAVETARNAS HÄRJARVISA}\\
Nu ska vi kompilera,\\
supa och penetrera\\
söka rekusioner ända in\\
till fixpunktssemantik.\\
Av induktion och basfall\\
blir man med lätthet asknall.\\
Hashtabeller parsas rekursivt\\
ifrån hår till häl\par
\vspace{10pt}
{\footnotesize\textit{Text: Pär Mattsson och Magnus Ingelbo från DVL
    Uppsala.}}


%% (Med handen för hjärtat står man upp och
%%     fortsätter med)}}
%% \vspace{10pt}
%% \fvisa{DATAVETARNAS HÄRJARVISA}{Nu ska vi kompilera}
%% \vspace{10pt}
%% Nu ska vi kompilera,\\
%% supa och penetrera\\
%% söka rekusioner ända in\\
%% till fixpunktssemantik.\\
%% Av induktion och basfall\\
%% blir man med lätthet asknall.\\
%% Hashtabeller parsas rekursivt\\
%% ifrån hår till häl\\
%% \\
%%   {\footnotesize\textit{Text: Pär Mattsson och Magnus Ingelbo från DVL Uppsala.}}
  
\chapterpage{Säsongsvisor}{Säsongsvisor}

\chapterpage{Svenska visor}{Svenska visor}

\chapterpage{Osvenska visor}{Osvenska visor}

\chapterpage{Phulavisor}{Phulavisor}

\chapterpage{Register}{Register}

\cfoot{}
%\newgeometry{top=1cm,bottom=1.2cm,left=1cm,right=1cm,footskip=0pt,headsep=0.2cm}
\small
\printindex
\hiddenchapterpage{$\emptyset$}
\ohead{}
\newpage
\null
\newpage
\null
\newpage
\null
\newpage
\null
\newpage
\null
\newpage
\null
\newpage
\null
\newpage
\null
\end{document}
